% ============================================================================
\section{The method in detail}
\label{app:method}

This appendix gives the details of the construction of \cref{sec:method} and illustrates its steps on synthetic swarms;
\cref{fig:overview} follows the same steps on the robot data.

\subsection{The footprint and the moments}

The distribution of the agents is a set of points, which the footprint turns into continuous fields. Four properties
are asked of it. It must favor no direction, so that the fields do not depend on the orientation of the scene. It must
be smooth, so that the fields can be differentiated with respect to the positions of the agents. It must be negligible
beyond a finite window, so that a field is evaluated locally. And the potentials of the flux must follow from it in
closed form, by solving a Poisson equation without a grid. A Gaussian footprint has all four
(\cref{prop:hhd,prop:so2,prop:tail}). The same philosophy, observing a signal at all scales at once through Gaussians of
growing width, underlies the scale space of image analysis \citep{witkin1983scale,lindeberg1994scale}; here, however,
the Gaussian was chosen for the four properties above, and its width as a lens is only one of its roles. Up to a
constant factor, $M_0$ and $\vM_1$ are the density and the momentum density by which discrete particle systems are
coarse-grained into continua, where a Gaussian is also chosen for its smooth fields and exact integrals
\citep{weinhart2012discrete}.

Two footprints whose centers are $d$ apart are no longer resolved, in the sense that their sum at the midpoint exceeds
the height of a single footprint, when $d\le2\sqrt{2\ln2}\,\sigma\approx2.35\,\sigma$ (\cref{fig:footprint}). At the
smallest width of the ladder, $d_{\min}/2$, two agents are therefore resolved as soon as they are more than
$\sqrt{2\ln2}\,d_{\min}\approx1.18\,d_{\min}$ apart. The floor $\epsilon_\sigma=0.01$\,m of \eqref{eq:ladder} is the
tracking precision.

\begin{figure}[!htbp]
\centering
\figin{\linewidth}{figures/fig_footprint.pdf}
\caption{The Gaussian footprint: (a)~one agent, evaluated within $5\sigma$; (b)~its profile; (c)~two agents at
distance $d$, no longer resolved below $2.35\,\sigma$; (d,\,e)~the width as a lens, from the agents to the group.}
\label{fig:footprint}
\end{figure}
\FloatBarrier

Each moment captures something that the others miss (\cref{fig:moments}). A group at rest has $\vM_1=\vec 0$ but
occupies space in $M_0$. Two streams crossing in opposite directions cancel in $\vM_1$ but add up in
$\operatorname{tr}\tM_2=\sum_i\|\vv_i\|^2O_i$, the local kinetic intensity. Likewise, a concentrated group and a
dispersed one can share the same net flux and still differ in $M_0$. The traceless part of $\tM_2$, which measures how
strongly the velocities align along a common axis, is also invariant under rotations, but we leave it aside.

\begin{figure}[!htbp]
\centering
\figin{\linewidth}{figures/fig_moments.pdf}
\caption{The three moments on a synthetic configuration at one instant ($\sigma=6.5$\,cm): a group at rest, two
crossing streams and a group moving together (a), and $M_0$, $\|\vM_1\|$ and $\operatorname{tr}\tM_2$ (b--d).}
\label{fig:moments}
\end{figure}
\FloatBarrier

\subsection{The two potentials}

With the sign of \eqref{eq:hhd}, the flux points toward the maxima of the divergence potential $\varphi_\sigma$, which
mark sinks where the agents converge, and away from its minima, which mark sources. The extrema of the vorticity
potential $\psi_\sigma$ are centers of rotation: minima for a counterclockwise rotation and maxima for a clockwise one.
On a bounded domain, the Helmholtz--Hodge decomposition is not unique. It is fixed by boundary conditions, or by setting
aside a harmonic part that stands for the influences from outside the domain \citep{bhatia2014natural}. Here the flux
is a sum of Gaussians, smooth on the whole plane and vanishing far from the agents. Such a field has no harmonic part
\citep[Section~3.1]{bhatia2013helmholtz}, and its two potentials are the solutions that vanish at infinity of the
Poisson equations $\Delta\varphi_\sigma=\nabla\cdot\vVb_\sigma$ and $\Delta\psi_\sigma=\nabla^{\perp}\cdot\vVb_\sigma$
\citep[Eq.~16]{bhatia2013helmholtz}, where $\nabla^{\perp}\cdot\vW=\partial_xW_y-\partial_yW_x$ is the two-dimensional
curl of a vector field $\vW$. The decomposition \eqref{eq:hhd} is then the two-dimensional form of the decomposition
without its harmonic part \citep[Eq.~15]{bhatia2013helmholtz}. A Gaussian is, up to a factor, the Laplacian of a radial
function whose derivative has a closed form, which yields the kernel $\kappa_\sigma$ of \eqref{eq:kappa}
(\cref{app:proofs}); its numerator cancels the singularity of $1/\|\vd\|^2$, so that it equals $1/2$ at $\vd=\vec 0$.
The two potentials \eqref{eq:potentials} rebuild the averaged flux exactly, with no harmonic remainder
(\cref{prop:hhd}). The four fields have no unit, and a replay of the recording at another constant rate, which
multiplies every velocity by the same factor, leaves them unchanged.

The potentials of a Helmholtz--Hodge decomposition also serve to locate the singularities of a vector field
\citep{polthier2003identifying}. They are usually computed on meshes \citep{polthier2003identifying,tong2003discrete},
with particles \citep{petronetto2010meshless} or with matrix-valued kernels fitted to scattered samples
\citep{macedo2010learning}, as reviewed by \citet{bhatia2013helmholtz}; the Gaussian footprint gives them in closed form
instead. \Cref{fig:flows} shows them on three elementary motions.

\begin{figure}[!htbp]
\centering
\figin{\linewidth}{figures/fig_flows.pdf}
\caption{The divergence and vorticity potentials of three elementary motions of $20$ synthetic agents ($\sigma=7$\,cm;
arrows: the averaged flux): a sink is a maximum of $\varphi_\sigma$, a source a minimum, and a counterclockwise
rotation center a minimum of $\psi_\sigma$. The potential that a motion does not produce stays below
$\nm{flows-flat-pct}\,\%$ of the color scale of its row.}
\label{fig:flows}
\end{figure}
\FloatBarrier

\subsection{The coupling}

Taken separately, the diagrams of $\varphi_\sigma$ and $\psi_\sigma$ record how strong a sink and a vortex are, but not
whether the two lie at the same place. Multiparameter persistence would summarize several fields together, but its
summaries do not reduce to a diagram, and comparing them is far less simple
\citep{carlsson2009theory,botnan2023introduction}. The coupling $J_\sigma$ of \eqref{eq:jacobi} is not a new physical
quantity, but it records in one field what the separate diagrams miss. Like the other fields, it is invariant under
rotations and translations, since a rotation does not change a determinant. \Cref{fig:potentials} shows two swarms with
the same ingredients, arranged differently: the diagrams of $\varphi_\sigma$ and $\psi_\sigma$ change only in part, while
that of $J_\sigma$ is empty when the sink and the rotation center coincide and appears when they are apart.


\subsection{Extended persistence and the sliced distance}

Persistent homology follows how the sublevel sets $\{f\leq c\}$ of a field change as the level $c$ rises
(\cref{fig:persistence}). A component is born at each minimum and dies when it merges with an older one. Cycles are born
at saddles and filled at maxima. Ordinary persistence \citep{edelsbrunner2002topological,zomorodian2005computing}
records the level at which each homology class, a connected component or a cycle, is born and the level at which it
dies. Some classes never die, such as the component that holds the global minimum; they are called essential, and their
death would be infinite. Extended persistence \citep{cohen2009extending} completes the ascending sequence of sublevel
sets with the homology of the domain relative to its superlevel sets, swept downward, so that every class has a finite
birth and death. It yields six diagrams per field: the ordinary diagrams $\mathrm{Ord}_0$ and $\mathrm{Ord}_1$, the
relative diagrams $\mathrm{Rel}_1$ and $\mathrm{Rel}_2$, and the extended diagrams $\mathrm{Ext}_0$ and
$\mathrm{Ext}_1$. On a fixed complex, the bottleneck distance between the extended diagrams of two fields is at most the
largest difference of their values \citep[Section~6]{cohen2009extending}. Persistence is multiscale in another sense
than the ladder: at one width, it follows the field across all its levels.

The sliced Wasserstein distance \eqref{eq:sw} is of order one. A field holds a few long-lived structures and many
short-lived ones, whose number and levels also describe the motion, and squared differences would let a single
long-lived structure outweigh all of them. The distance is bounded above by $2\sqrt2$ times the first diagram distance
\citep[Theorem~3.3]{carriere2017sliced}, the least total displacement that matches the points of two diagrams with one
another or with the diagonal, so that a small displacement of the points changes it little. It is differentiable
almost everywhere, since the sort and the absolute value are linear by pieces.

\begin{figure}[!htbp]
\centering
\figin{\linewidth}{figures/fig_persistence.pdf}
\caption{From a field to a distance: (a)~the sublevel sets of a profile; (b)~the resulting points, ordinary and
extended; (c)~two diagrams, each augmented with the projections of the other onto the diagonal (hollow); (d)~their
sorted projections on one direction.}
\label{fig:persistence}
\end{figure}
\FloatBarrier

\subsection{The domain without a grid}

Each recording gets its own seeds at every width, of three kinds. The first kind is a subset of the agent positions,
chosen by farthest-point sampling until every position lies within $\nm{cover-sigma}$ of the subset. The second, for
each field, are the barycenters of the $\nm{frac-pct}\%$ of the Delaunay triangles of these positions where its
gradient is smallest; these seeds lie near the critical points of the field. The third are up to $\nm{n-jacobi}$ points
of the coupling skeleton, located by $\nm{n-bis}$ bisection steps on the Delaunay edges where $J_\sigma$ changes sign.
The five fields are evaluated at the same seeds, and seeds farther than $5\sigma_m$ from every agent at every instant
are discarded. The domain is the alpha complex of the remaining seeds, computed with GUDHI \citep{maria2014gudhi}. It
keeps the simplices of their Delaunay triangulation whose alpha value, that is, whose squared radius, is at most
$(\nm{c-alpha}\,\sigma_m)^2$ \citep{edelsbrunner2010computational}. On this complex, each field $f$ induces a lower-star
filtration: a vertex enters at the value of $f$ there, and an edge or a triangle at the largest value of $f$ at its
vertices.

\subsection{Ranks, fusion, affinity propagation and the modal region}

With mutual ranks, a pair of recordings is close only when each is among the nearest of the other, and the percentile
ranks then give every matrix the same distribution. Weights read from the data, through the spectral gap of each field,
are compared with the weights of \eqref{eq:view} in \cref{app:fusion}. In the fusion \eqref{eq:snf}, the local kernel
keeps distinct behaviors apart, while the diffusion joins the recordings of one behavior that single widths leave split.

Affinity propagation passes two kinds of messages, responsibilities and availabilities, between the recordings. The
messages are updated until a set of exemplars emerges and each recording is assigned to one of them. Its preference,
which sets how readily a recording becomes an exemplar, is the median of the entries of $\Sstar$, its diagonal first set
to its largest value. Its damping and its convergence rule are those of its original implementation
(\cref{tab:settings}), so that the clustering adds no tuned parameter. In the modal consensus, ties for the
exemplar of a group go to the member most often clustered with the others.

\subsection{Splitting and differentiating a distance}

Each term of the sliced distance is shared equally between the two elements it matches, a projection onto the diagonal
standing for the point it projects. Each diagram point then passes its share, half and half, to the two seeds that
realize its birth and its death. Holding the seeds fixed while the trajectories move is exact at first order for the
critical values, whose gradient vanishes. It is not exact for a quantile of the values, which lies at an ordinary point
where the gradient does not vanish (\cref{app:attribution}).
